\begin{tabular}{p{2.6cm}cp{1.5cm}p{8.4cm}}
\toprule
Study-arm & $n$ & Follow-up & Clinical-success definition as recorded from the source \\
\midrule
Maddocks2019\_A & 1 & 3 d & Improvement of oxygenation within 3 days with successful ventilator weaning, in a 77-year-old woman with VAP and empyema complicated by bronchopleural fistula after thoracotomy \\
Blasco2023\_A & 1 & 7 d & Treatment FAILURE (paper title): new P. aeruginosa BSI with septic residual limb metastasis occurred in absence of antimicrobial therapy after the 1-week phage course; ultimately required vascular prosthesis replacement \\
Ngauy2026\_A & 1 & 2 wk & Clinical and microbiological resolution of hypermutator XDR osteomyelitis/hardware infection; ESR/CRP normalized after 1 week; orbital implant culture-negative at removal 2 weeks post-phage \\
Yang2025\_A & 1 & 2 wk & Author-reported significant clinical improvement after 2 weeks of nebulized phage therapy (vasopressors discontinued, haemodynamics stable, inflammatory markers and lung effusion decreased). \\
Aslam2020\_C10 & 1 & 3 mo & Recurrent P. aeruginosa bacteremia from a probable aortic graft infection; blood cultures negative on therapy and four weekly surveillance cultures negative with no recurrence through 12 weeks after end of therapy, against a pre-treatment pattern of recurrence 7-10 days off antibiotics \\
Teney2024\_A & 1 & 7 mo & 52-year-old man with self-inflicted burns over 81\% of total body surface area and four successive episodes of ventilator-associated pneumonia with bacteraemia from an NDM-1-producing XDR strain. Success recorded as reduced bronchorrhoea and oxygen dependence, improving skin grafts, extubation on the last day of the second phage course, and discharge from intensive care after 203 days \\
Racenis2022\_femur\_A & 1 & 9 mo & Eradication of infection in PROXIMAL femoral segment/hip region (site treated with local phage BFC 1.10); wound healed with no local/systemic infection signs at days 1-15 post-phage; enabled endoprosthesis reconstruction 9 months later. IMPORTANT: a SEPARATE, untreated-by-phage recurrence of MDR P. aeruginosa+VRE occurred at the DISTAL femoral segment ~9 months later and was treated with antibiotics alone (DAIR, colistin+fosfomycin) - that distal episode is NOT part of this phage-treated arm and is excluded from this row's outcome counts \\
Law2019\_A & 1 & 9 mo & Author-reported clinical resolution of MDR P. aeruginosa pneumonia; no recurrence of pseudomonal pneumonia or CF exacerbation within 100 days of ending bacteriophage therapy; successful bilateral lung transplantation 9 months later. \\
Tkhilaishvili2020\_A & 1 & 10 mo & Infection eradicated; no pain, satisfactory mobility, normal CRP, no implant loosening on X-ray at 10-month follow-up \\
Ferry2021\_A & 1 & 12 mo & Resolution of relapsing prosthetic knee infection: normal local knee status, painless motion and walking at 1-year follow-up, on suppressive antimicrobial therapy \\
Chan2018\_omko1\_A & 1 & 18 mo & Chronic MDR P. aeruginosa aortic Dacron graft infection resolved after a single local application of phage OMKO1 + ceftazidime; no evidence of recurrent infection over 18 months off all antibiotics (blood cultures, clinical symptoms, CT). \\
Denis2026\_A & 1 & 2.0 yr & Author-reported cure: at 2-year follow-up the patient had no relapse of the previously thrice-relapsing MDR P. aeruginosa frontal osteitis, without any suppressive treatment. \\
Arya2026\_pji\_A & 1 & 2.0 yr & Author-reported clinical resolution of a multidrug-resistant, non-operable P. aeruginosa periprosthetic joint infection that had failed multiple antibiotic and surgical interventions, achieved with intermittent bacteriophage therapy alone (no antibiotics) over ~2 years; microbiological eradication was NOT achieved (chronic suppression, altered virulence, biofilm disruption). \\
Cesta2023\_A & 1 & 2.0 yr & Chronic hip prosthetic joint infection: at 2-year follow-up after suspension of therapy there were no clinical signs of relapse and a labelled-leukocyte scan showed no pathological uptake \\
Leveque2023\_A & 1 & 2.0 yr & Clinical failure: despite microbiological eradication of the XDR strain, the patient progressed to multi-organ failure and died at 24 months post-lung-transplant; the chronic bronchopulmonary infection favored the fatal course. \\
Liu2025\_perinephric\_P2 & 1 & not stated & Treatment FAILURE for primary abscess site: urinary tract infection resolved and inflammatory area reduced (72x15mm to 21x10mm on ultrasound) but perinephric P. aeruginosa PERSISTED in wound secretions despite 3 courses of phage therapy; patient discharged per personal choice, not cured \\
Onallah2023\_PASA16\_agg & 15 & not stated & 'Good clinical outcome' as defined by study authors (binary); 2 clinical failures among 15 \\
Aslam2019\_P1 & 1 & not stated & Clinical improvement: discharged from hospital, tracheostomy de-cannulated, successful rehabilitation \\
Aslam2019\_P2 & 1 & not stated & Clinical improvement: weaned from ventilator by day 90, discharged from hospital \\
Pirnay2024\_XDR & 7 & not stated & 'Clinical improvement' as defined by study authors (binary yes/no); patients (case \#) 8,13,18,24,26,79 improved, patients 3 did not \\
Pirnay2024\_MDR & 23 & not stated & 'Clinical improvement' as defined by study authors (binary yes/no); patients (case \#) 4,11,17,20,22,27,34,54,56,64,67,75,81,82,91,92 improved, patients 14,36,38,76,87,93,96 did not \\
Pirnay2024\_PDR & 2 & not stated & 'Clinical improvement' as defined by study authors (binary yes/no); patients (case \#) 74 improved, patients 57 did not \\
Kohler2023\_A & 1 & not stated & Author-reported clinical improvement (binary yes/no): drastic improvement in general condition, progressive radiological clearance of lung consolidations on serial CT, resolution of respiratory exacerbations, and eventual complete discontinuation of systemic antibiotics; patient alive and returned to work at last follow-up (D996, Dec 2022). \\
Chan2025\_MDR & 7 & not stated & Not extractable as a binary clinical-success proportion: the study's endpoints are continuous (percent-predicted FEV1 change and sputum P. aeruginosa CFU reduction), reported cohort-wide, with no binary clinical-cure/success definition -- coded NA on the same basis as the BX004-A and PhagoBurn trials. \\
Chan2025\_PDR & 2 & not stated & Not extractable as a binary clinical-success proportion: the study's endpoints are continuous (percent-predicted FEV1 change and sputum P. aeruginosa CFU reduction), reported cohort-wide, with no binary clinical-cure/success definition -- coded NA on the same basis as the BX004-A and PhagoBurn trials. \\
Hahn2023\_A & 2 & not stated & Not extractable as a binary clinical-success proportion: the case series reports safety and radiological/CFU trends for two CF patients without a clean per-patient binary cure/success statement. \\
Duplessis2018\_A & 1 & not stated & The bacteremia was sterilized by phage, but the overall outcome was fatal (Fatal Outcome per the source); clinical success coded 0 given death despite microbiological clearance. \\
Malhotra2026\_A & 1 & not stated & Author-reported successful use of bacteriophage therapy for a recalcitrant MDR P. aeruginosa intra-abdominal infection after multiorgan transplant. \\
Li2025\_biliary\_A & 1 & not stated & Author-reported marked clinical improvement: abdominal pain alleviated, bile color normalized, and inflammatory markers (CRP, PCT) decreased after personalized P. aeruginosa phage therapy plus antibiotics for a complex recurrent biliary tract infection. \\
Aslam2020\_C5 & 1 & not stated & Ventricular assist device driveline infection with sternal osteomyelitis and recurrent bacteremia; explicitly recorded as Outcome: Failure -- bacteremia recurred one week after starting phage therapy and driveline drainage recurred after therapy ended \\
Aslam2020\_C8 & 1 & not stated & Persistently infected ventricular assist device with recurrent bacteremia; explicitly recorded as Outcome: Failure -- bacteremia recurred within one week of ending the first phage course and again four weeks into the second course while still on therapy \\
Khatami2021\_A & 1 & not stated & Extensive chronic Pseudomonas aeruginosa osteoarticular infection in a 7-year-old: a significant therapeutic effect was documented within two weeks of starting adjunctive phage therapy, averting the catastrophic outcome the therapy was attempted to avoid \\
Zaldastanishvili2021\_P1 & 1 & not stated & 43-year-old male with cystic fibrosis and chronic P. aeruginosa lower respiratory tract infection managed with IV antibiotics since childhood. Success is recorded as replacement of antibiotics by phages with subjective symptom alleviation and a 10- to 100-fold reduction in sputum bacterial load; the organism was never absent from any sputum culture over 2017-2020 \\
Zaldastanishvili2021\_P2 & 1 & not stated & 64-year-old female with primary ciliary dyskinesia and bronchiectasis who refused antibiotic therapy options. Success is recorded as management of the infection with phages alone across five successive custom preparations (2018-2021) with subjective symptom alleviation and no adverse effects \\
Green2023\_C11 & 1 & not stated & Persistent disseminated left-ventricular-assist-device driveline infection (strain MYC4). Recorded by the source as clinical improvement: computed tomography shows no evidence of infection and the patient has been off antibiotic treatment since May 2022 and stable, while scant driveline drainage remains culture-positive and is controlled with local wound care \\
Casazza2025\_A & 1 & not stated & 47-year-old woman with chronic mastoiditis from MDR P. aeruginosa on a background of cystic-fibrosis ciliary dysfunction and lung-transplant immunosuppression. Success recorded as resolution of infection confirmed by symptomatology, cultures and imaging: at five months, resolution of otorrhoea, headache and hearing impairment \\
Li2023\_ILD\_A & 1 & not stated & 40-year-old man with progressive interstitial lung disease, type I respiratory failure and chronic infection by a pyomelanin-producing MDR strain resistant to all 100 dsDNA phages in the treating library. Success recorded as decreased cough, expectoration and dyspnoea, reduced bacterial burden, complete cessation of antibiotics, and discharge with relieved symptoms; lung transplantation six months later with no lung infection on follow-up \\
Chen2022\_empyema\_A & 1 & not stated & 68-year-old man with post-pneumonectomy empyema and broncho-pleural fistula from carbapenem-resistant P. aeruginosa. Success recorded as clearance of the pathogen and improvement of the clinical outcome \\
Rubalskii2020\_P8 & 1 & not stated & 13-year-old male with a P. aeruginosa-infected thoracotomy wound two months after double lung transplantation for cystic fibrosis, refractory to surgical debridement, vacuum-assisted closure and continuous antibiotic therapy under mycophenolic acid, tacrolimus and prednisolone. Success recorded as the wound healing completely \\
PardoFreire2025\_A & 1 & not stated & 44-year-old woman with cystic fibrosis and bilateral lung transplantation, in acute rejection with deteriorating lung function on a background of chronic P. aeruginosa. Success recorded as mucus clearance, improved pulmonary function and resolution of the acute rejection, sustained to one year \\
\bottomrule
\end{tabular}
